\scanpage{65}
% archive sheet number: 67

[On laisse tomber dorénavant les $'$ de $X'$!]

\deleted{Mais sur $X^*=X\setminus\bigcup D_i$} (si \deleted{\unsure{disques compacts|chaque composante} \unsure{connexes}} \inserted{OPS} $X$ connexe, $I$ non vide --- alors \deleted{\unsure{$X$ n'est}} il y a sur $X^*=X\setminus\bigcup_i\mathring D_i$ (surface conforme à bord \unsure{contenue|continue|contenu} dans la découpe de $X$, suivant $\bigcup\partial D_i$) une \deleted{\unsure{uniq}} métrique \deleted{hyperbolique} canonique à courbure constante (qu'on peut d'ailleurs normaliser par la condition que les volumes sont $1$) \deleted{\unsure{avec des|sont des}} \inserted{et $=$} canonique \inserted{sans cts} dans le cas hyperbolique (i.e. quand on n'a \deleted{\unsure{pas}} $g=0\quad\underset{\mathrm{card}\,I}{\underset{\scriptscriptstyle\parallel}{\nu}}\in\{1,2\}$), dans le cas hyperbolique le\deleted{\unsure{s}} bord est géodésique,
% slanted
\margin{il y a aussi le cas \deleted{\unsure{d'un|deux}} \unsure{oublié} du plan projectif réel $(\simeq S^2/\pm)$ \deleted{\unsure{privé} d'un disque} ``avec un trou''}
dans tous les cas le bord est muni d'une structure torique canonique provenant de la structure riemannienne.

Il est tentant de restreindre le\inserted{ syst\textsuperscript{e} des} disques envisagé, en exigeant que les deux structures toriques sur \deleted{\unsure{à}} chaque $\partial D_i$ (provenant soit de $D_i$, soit de $X^*$)
